\documentclass[tikz,border=4pt]{standalone}
\usepackage{ctex}
\usepackage{amsmath}
\usetikzlibrary{shapes.geometric, arrows.meta, positioning}

\tikzset{
  box/.style={rectangle, draw=black, rounded corners=3pt,
              minimum width=3.8cm, minimum height=0.65cm,
              align=center, font=\small},
  diam/.style={diamond, draw=black, aspect=2.0,
               minimum width=4.5cm, minimum height=1.0cm,
               align=center, font=\small},
  arr/.style={->, >=Stealth, thick},
}

\begin{document}
% 直接开平方法与因式分解法选择流程
\begin{tikzpicture}[node distance=1.0cm]

  \node[box, fill=blue!10, minimum width=7cm] (start)
    {$ax^2 + bx + c = 0$，先整理为标准形式};

  \node[diam, below=1.2cm of start, fill=yellow!20] (q1)
    {形如 $(x-p)^2=q$ 或 $ax^2=k$？};

  \node[box, fill=green!15, below left=1.2cm and 1.8cm of q1] (sqrt)
    {\textbf{直接开平方法}\\$(x-p)^2=q \Rightarrow x-p=\pm\sqrt{q}$};

  \node[diam, below right=1.2cm and 1.0cm of q1, fill=yellow!20] (q2)
    {整数系数能十字相乘？};

  \node[box, fill=green!15, below left=1.0cm and 0.5cm of q2] (factor)
    {\textbf{因式分解法}\\$(x-x_1)(x-x_2)=0$};

  \node[box, fill=orange!20, below right=1.0cm and 0.5cm of q2] (formula)
    {\textbf{配方法/求根公式}\\（通用）};

  \draw[arr] (start) -- (q1);
  \draw[arr] (q1) -- node[above left,font=\small]{是} (sqrt);
  \draw[arr] (q1) -- node[above right,font=\small]{否} (q2);
  \draw[arr] (q2) -- node[above left,font=\small]{是} (factor);
  \draw[arr] (q2) -- node[above right,font=\small]{否} (formula);

\end{tikzpicture}
\end{document}
